\documentclass[11pt]{article}
\usepackage[utf8]{inputenc}
\usepackage[IL2]{fontenc}
\usepackage[czech]{babel}
\usepackage{a4wide}
\usepackage{amsthm,amsmath,amsfonts,amssymb}
\usepackage{xcolor}
\usepackage{soul}
\usepackage{graphicx}
\title{DÚ LA 5.1}
\date{}
\begin{document}
\maketitle
\paragraph{} Podmínku na součet prvků v každém sloupci matice $A$ zapíšeme jako maticové násobení a upravíme:
\begin{align*}
A^T\begin{pmatrix}
1\\1\\\vdots\\1
\end{pmatrix}=\begin{pmatrix}
a\\a\\\vdots\\a
\end{pmatrix}=&\begin{pmatrix}
  a      & 0      & \dots  & 0      \\
  0      & a      & \dots  & 0      \\
  \vdots & \vdots & \ddots & \vdots \\
  0      & 0      & \dots  & a
\end{pmatrix}\begin{pmatrix}
1\\1\\\vdots\\1
\end{pmatrix}\\
\left(A^T-\begin{pmatrix}
  a      & 0      & \dots  & 0      \\
  0      & a      & \dots  & 0      \\
  \vdots & \vdots & \ddots & \vdots \\
  0      & 0      & \dots  & a
\end{pmatrix}\right)\begin{pmatrix}
1\\1\\\vdots\\1
\end{pmatrix}=&\textbf{o}
\end{align*}
\paragraph*{} Označme
\[
B^T:=A^T-\begin{pmatrix}
  a      & 0      & \dots  & 0      \\
  0      & a      & \dots  & 0      \\
  \vdots & \vdots & \ddots & \vdots \\
  0      & 0      & \dots  & a
\end{pmatrix}
\]
\paragraph*{} Má tedy platit, že $B^T(1, 1,\dots,1)^T=\textbf{o}$, tj. že 0 je vlastním číslem matice $B^T$, a tedy (dle pozorování 9.14) je $B^T$ singulární. Proto je singulární i $B$,\footnote{Pro libovolnou čtvercovou matici $X$ můžeme rovnost $XX^{-1}=I_n$ upravit transponováním na $(X^{-1})^TX^T=I_n$, a tedy $X$ je regulární, právě když je $X^T$ regulární.} a tedy $B$ má vlastní číslo 0 (opět pozorování 9.14), tj. $\exists\;\textbf{o}\neq\textbf{x}\in\textbf{T}^n:B\textbf{x}=\textbf{o}$. Dle tvrzení 4.16, bodu (2) můžeme $B$ nahradit rozdílem $A-$diag$(a, a, \dots, a)$ a dále upravovat:
\begin{align*}
\left(A-\begin{pmatrix}
  a      & 0      & \dots  & 0      \\
  0      & a      & \dots  & 0      \\
  \vdots & \vdots & \ddots & \vdots \\
  0      & 0      & \dots  & a
\end{pmatrix}\right)\textbf{x}&=\textbf{o}\\
A\textbf{x}-\begin{pmatrix}
  a      & 0      & \dots  & 0      \\
  0      & a      & \dots  & 0      \\
  \vdots & \vdots & \ddots & \vdots \\
  0      & 0      & \dots  & a
\end{pmatrix}\textbf{x}&=\textbf{o}\end{align*}
\begin{align*}
A\textbf{x}&=\begin{pmatrix}
  a      & 0      & \dots  & 0      \\
  0      & a      & \dots  & 0      \\
  \vdots & \vdots & \ddots & \vdots \\
  0      & 0      & \dots  & a
\end{pmatrix}\textbf{x}\\
A\textbf{x}&=a\textbf{x}
\end{align*}
\paragraph*{} A protože $\textbf{x}\neq\textbf{o}$, je $a$ vlastní číslo matice $A$, což jsme chtěli dokázat.\qed
\end{document}